\chapter{Complexity — Path Walk Machine}
%%% generated by the blueprint pipeline from TCSlib.Complexity.SpaceComplexity.Logspace.PathWalkMachine
\section{Overview}

Declarations of the Lean module \texttt{TCSlib.Complexity.SpaceComplexity.Logspace.PathWalkMachine}.

\section{Declarations}

\begin{definition}[Path program labels]
\lean{Complexity.LogProg.PathLabel}
\label{Complexity.LogProg.PathLabel}
\leanok
Introduce a finite set of labels consisting of the listed unary-block, matrix-block,
source, target, walk, acceptance, and rejection labels, together with four families of
labels indexed by $k\in\{0,1,2\}$.
\end{definition}

\begin{definition}[Three-entry index map]
\lean{Complexity.LogProg.akY}
\label{Complexity.LogProg.akY}
\leanok
Define a map from $\{0,1,2\}$ to $\{0,1,\ldots,8\}$ whose values at $0,1,2$ are
respectively $1,1,8$.
\end{definition}

\begin{definition}[Example label map]
\lean{Complexity.LogProg.akEx}
\label{Complexity.LogProg.akEx}
\leanok
\uses{Complexity.LogProg.PathLabel}
Define a map from the three-element set $	ext{Fin }3$ to the path labels by sending its
three elements, in order, to $wi1$, $wrn$, and $wrd$.
\end{definition}

\begin{definition}[Nine-register program]
\lean{Complexity.LogProg.pathA}
\label{Complexity.LogProg.pathA}
\leanok
\uses{Complexity.LogProg.NARM, Complexity.LogProg.PathLabel, Complexity.LogProg.akEx, Complexity.LogProg.akY}
Define a nondeterministic register program with nine natural-number registers and two
available deciders, whose control locations are the path labels. At each location its
instruction is the one specified by the defining clauses: these clauses give the
register operations, branches, decider calls, choices, and halt answers.
\end{definition}

\begin{lemma}[Path program precondition]
\lean{Complexity.LogProg.pathA_npreS}
\label{Complexity.LogProg.pathA_npreS}
\leanok
\uses{Complexity.LogProg.AConf, Complexity.LogProg.NARM.detPart, Complexity.LogProg.NPreS, Complexity.LogProg.PathLabel, Complexity.LogProg.PreS, Complexity.LogProg.akY, Complexity.LogProg.pathA}
\difficulty{5}
For every bit string $x$ and every configuration consisting of an optional control
label, nine natural-number registers, and an optional choice bit, the precondition for
the program obtained from `pathA` by replacing each choice instruction with the answer
“no” holds. In particular, every register comparison used by the program compares
distinct registers, every oracle call has no repeated register argument, and any
input-format check required at the current instruction is satisfied.
\end{lemma}

\begin{definition}[Input languages]
\lean{Complexity.LogProg.pathAs}
\label{Complexity.LogProg.pathAs}
\leanok
\uses{Complexity.inputBitLang, Complexity.inputLenLang}
Define a map from the two-element set $	ext{Fin }2$ to languages of Boolean words by
assigning the input-bit language to the first element and the input-length language to
the second.
\end{definition}

\begin{definition}[Language indicator map]
\lean{Complexity.LogProg.pathOr}
\label{Complexity.LogProg.pathOr}
\leanok
\uses{Complexity.LogProg.pathAs, Turing.MultiTapeTM.indicator}
For each of the two languages just specified and each Boolean word $V$, define its
indicator to be $1$ when $V$ belongs to that language and $0$ otherwise.
\end{definition}

\begin{lemma}[Indexed bit test]
\lean{Complexity.LogProg.pathOr_bit}
\label{Complexity.LogProg.pathOr_bit}
\leanok
\uses{Complexity.LogProg.pathAs, Complexity.LogProg.pathOr, Complexity.indicator_indexLang_pairEncode, Complexity.inputBitLang, Turing.pairEncode}
\difficulty{1}
For every bit string $x$ and natural number $i$, the Boolean indicator of the language
that tests whether the indexed bit is $1$, evaluated on the encoding of $x$ and the
binary representation of $i$, equals the bit of $x$ at index $i$, with false returned
when that index is beyond the end of $x$.
\end{lemma}

\begin{lemma}[Indexed length test]
\lean{Complexity.LogProg.pathOr_len}
\label{Complexity.LogProg.pathOr_len}
\leanok
\uses{Complexity.LogProg.pathAs, Complexity.LogProg.pathOr, Complexity.indicator_indexLang_pairEncode, Complexity.inputLenLang, Turing.pairEncode}
\difficulty{1}
For every bit string $x$ and natural number $i$, the Boolean indicator of the language
that tests whether the index is below the input length, evaluated on the encoding of $x$
and the binary representation of $i$, is true exactly when $i<|x|$.
\end{lemma}

\begin{lemma}[Log-space membership of the path languages]
\lean{Complexity.LogProg.pathAs_mem}
\label{Complexity.LogProg.pathAs_mem}
\leanok
\uses{Complexity.LOGSPACE, Complexity.LogProg.pathAs, Complexity.inputBitLang_mem, Complexity.inputLenLang_mem}
\difficulty{4}
For each $j\in\{0,1\}$, the language that, on an input encoding a bit string $x$ and an
integer $i$, tests whether the $i$th bit of $x$ is $1$ when $j=0$, and whether $i<|x|$
when $j=1$, is decidable in logarithmic space up to a constant factor.
\end{lemma}

\begin{definition}[Program register state]
\lean{Complexity.LogProg.PReg}
\label{Complexity.LogProg.PReg}
\leanok
A program register state is a tuple of nine natural numbers $(p,n,a,b,s,t,tmp,st,nx)$,
representing respectively the input pointer, vertex count, row counter, column counter,
source or current vertex, target, scratch copy, step count, and guessed next vertex.
\end{definition}

\begin{definition}[Register value vector]
\lean{Complexity.LogProg.PReg.v}
\label{Complexity.LogProg.PReg.v}
\leanok
\uses{Complexity.LogProg.PReg}
For a program register state $\sigma$, define its value vector in $\mathbb{N}^9$ to list, in order,
its input pointer, vertex count, row counter, column counter, source or current vertex,
target, scratch copy, step count, and guessed next vertex.
\end{definition}

\begin{definition}[Register bound]
\lean{Complexity.LogProg.PReg.Le}
\label{Complexity.LogProg.PReg.Le}
\leanok
\uses{Complexity.LogProg.PReg}
For a program register state $\sigma$ and a natural number $B$, the state is bounded by $B$
exactly when each of its nine register values is at most $B$.
\end{definition}

\begin{lemma}[Register Bounds in a Configuration]
\lean{Complexity.LogProg.PReg.reg}
\label{Complexity.LogProg.PReg.reg}
\leanok
\uses{Complexity.LogProg.PReg, Complexity.LogProg.PReg.Le, Complexity.LogProg.PReg.v, Complexity.LogProg.PathLabel, Complexity.LogProg.RegLe}
\difficulty{4}
Let $\sigma$ be a nine-register state, let $B\in\mathbb{N}$, and suppose each register
of $\sigma$ is at most $B$. For any optional control label and optional Boolean result,
every register in the configuration formed from that label, the register vector of
$\sigma$, and the result is at most $B$.
\end{lemma}

\begin{lemma}[Input Pointer Register Value]
\lean{Complexity.LogProg.PReg.v_0}
\label{Complexity.LogProg.PReg.v_0}
\leanok
\uses{Complexity.LogProg.PReg.v}
\difficulty{1}
For every program register state, the value of its register vector at index $0$ is its
input pointer $p$.
\end{lemma}

\begin{lemma}[Vertex Count Register Value]
\lean{Complexity.LogProg.PReg.v_1}
\label{Complexity.LogProg.PReg.v_1}
\leanok
\uses{Complexity.LogProg.PReg.v}
\difficulty{1}
For every program register state, the value of its register vector at index $1$ is its
vertex count $n$.
\end{lemma}

\begin{lemma}[Row-counter register value]
\lean{Complexity.LogProg.PReg.v_2}
\label{Complexity.LogProg.PReg.v_2}
\leanok
\uses{Complexity.LogProg.PReg.v}
\difficulty{1}
For a program register state with row counter $a$, the value of its register-value
vector at index $2$ is $a$.
\end{lemma}

\begin{lemma}[Column-counter register value]
\lean{Complexity.LogProg.PReg.v_3}
\label{Complexity.LogProg.PReg.v_3}
\leanok
\uses{Complexity.LogProg.PReg.v}
\difficulty{1}
For a program register state with column counter $b$, the value of its register-value
vector at index $3$ is $b$.
\end{lemma}

\begin{lemma}[Current-vertex register value]
\lean{Complexity.LogProg.PReg.v_4}
\label{Complexity.LogProg.PReg.v_4}
\leanok
\uses{Complexity.LogProg.PReg.v}
\difficulty{1}
For a program register state with current-vertex value $s$, the value of its
register-value vector at index $4$ is $s$.
\end{lemma}

\begin{lemma}[Target register value]
\lean{Complexity.LogProg.PReg.v_5}
\label{Complexity.LogProg.PReg.v_5}
\leanok
\uses{Complexity.LogProg.PReg.v}
\difficulty{1}
For a program register state with target value $t$, the value of its register-value
vector at index $5$ is $t$.
\end{lemma}

\begin{lemma}[Scratch-copy register value]
\lean{Complexity.LogProg.PReg.v_6}
\label{Complexity.LogProg.PReg.v_6}
\leanok
\uses{Complexity.LogProg.PReg.v}
\difficulty{1}
For a program register state with scratch-copy value $tmp$, the value of its
register-value vector at index $6$ is $tmp$.
\end{lemma}

\begin{lemma}[Step-count register value]
\lean{Complexity.LogProg.PReg.v_7}
\label{Complexity.LogProg.PReg.v_7}
\leanok
\uses{Complexity.LogProg.PReg.v}
\difficulty{1}
For a program register state with step-count value $st$, the value of its register-value
vector at index $7$ is $st$.
\end{lemma}

\begin{lemma}[Guessed-next-vertex register value]
\lean{Complexity.LogProg.PReg.v_8}
\label{Complexity.LogProg.PReg.v_8}
\leanok
\uses{Complexity.LogProg.PReg.v}
\difficulty{1}
For a program register state with guessed-next-vertex value $nx$, the value of its
register-value vector at index $8$ is $nx$.
\end{lemma}

\begin{lemma}[Updating the Input Pointer]
\lean{Complexity.LogProg.PReg.u_0}
\label{Complexity.LogProg.PReg.u_0}
\leanok
\uses{Complexity.LogProg.PReg.v}
\difficulty{4}
Let $\sigma$ be a nine-register state and let $z\in\mathbb{N}$. Replacing coordinate $0$
of its register vector by $z$ gives the register vector of the state obtained by
replacing its input-pointer register by $z$.
\end{lemma}

\begin{lemma}[Updating the Vertex Count]
\lean{Complexity.LogProg.PReg.u_1}
\label{Complexity.LogProg.PReg.u_1}
\leanok
\uses{Complexity.LogProg.PReg.v}
\difficulty{4}
Let $\sigma$ be a nine-register state and let $z\in\mathbb{N}$. Replacing coordinate $1$
of its register vector by $z$ gives the register vector of the state obtained by
replacing its vertex-count register by $z$.
\end{lemma}

\begin{lemma}[Updating the Row Counter]
\lean{Complexity.LogProg.PReg.u_2}
\label{Complexity.LogProg.PReg.u_2}
\leanok
\uses{Complexity.LogProg.PReg.v}
\difficulty{4}
Let $\sigma$ be a nine-register state and let $z\in\mathbb{N}$. Replacing coordinate $2$
of its register vector by $z$ gives the register vector of the state obtained by
replacing its row-counter register by $z$.
\end{lemma}

\begin{lemma}[Updating the Column Counter]
\lean{Complexity.LogProg.PReg.u_3}
\label{Complexity.LogProg.PReg.u_3}
\leanok
\uses{Complexity.LogProg.PReg.v}
\difficulty{4}
Let $\sigma$ be a nine-register state and let $z\in\mathbb{N}$. Replacing coordinate $3$
of its register vector by $z$ gives the register vector of the state obtained by
replacing its column-counter register by $z$.
\end{lemma}

\begin{lemma}[Updating the Source Register]
\lean{Complexity.LogProg.PReg.u_4}
\label{Complexity.LogProg.PReg.u_4}
\leanok
\uses{Complexity.LogProg.PReg.v}
\difficulty{4}
For a program register state with value vector $(p,n,a,b,s,t,tmp,st,nx)$, replacing its
entry at index $4$ by $z\in\mathbb{N}$ gives the value vector of the state obtained by
replacing $s$ by $z$.
\end{lemma}

\begin{lemma}[Updating the Target Register]
\lean{Complexity.LogProg.PReg.u_5}
\label{Complexity.LogProg.PReg.u_5}
\leanok
\uses{Complexity.LogProg.PReg.v}
\difficulty{4}
For a program register state with value vector $(p,n,a,b,s,t,tmp,st,nx)$, replacing its
entry at index $5$ by $z\in\mathbb{N}$ gives the value vector of the state obtained by
replacing $t$ by $z$.
\end{lemma}

\begin{lemma}[Updating the Scratch Register]
\lean{Complexity.LogProg.PReg.u_6}
\label{Complexity.LogProg.PReg.u_6}
\leanok
\uses{Complexity.LogProg.PReg.v}
\difficulty{4}
For a program register state with value vector $(p,n,a,b,s,t,tmp,st,nx)$, replacing its
entry at index $6$ by $z\in\mathbb{N}$ gives the value vector of the state obtained by
replacing $tmp$ by $z$.
\end{lemma}

\begin{lemma}[Updating the Step Counter]
\lean{Complexity.LogProg.PReg.u_7}
\label{Complexity.LogProg.PReg.u_7}
\leanok
\uses{Complexity.LogProg.PReg.v}
\difficulty{4}
For a program register state with value vector $(p,n,a,b,s,t,tmp,st,nx)$, replacing its
entry at index $7$ by $z\in\mathbb{N}$ gives the value vector of the state obtained by
replacing $st$ by $z$.
\end{lemma}

\begin{lemma}[Updating the Guessed Vertex Register]
\lean{Complexity.LogProg.PReg.u_8}
\label{Complexity.LogProg.PReg.u_8}
\leanok
\uses{Complexity.LogProg.PReg.v}
\difficulty{4}
For a program register state with value vector $(p,n,a,b,s,t,tmp,st,nx)$, replacing its
entry at index $8$ by $z\in\mathbb{N}$ gives the value vector of the state obtained by
replacing $nx$ by $z$.
\end{lemma}

\begin{definition}[Program semantics predicate]
\lean{Complexity.LogProg.PS}
\label{Complexity.LogProg.PS}
\leanok
\uses{Complexity.LogProg.NSem, Complexity.LogProg.PReg, Complexity.LogProg.PReg.v, Complexity.LogProg.PathLabel, Complexity.LogProg.RegLe, Complexity.LogProg.pathA, Complexity.LogProg.pathOr}
Given a Boolean word $x$, a natural-number bound $B$, a path label $\ell$, a program
register state $\sigma$, and a proposition $P$, define the predicate to mean that $P$ holds
in the nondeterministic semantics of the specified nine-register program, starting at
label $\ell$ with register values $\sigma$ and no additional configuration data, under
register bound $B$.
\end{definition}

\begin{lemma}[Register Equality Preserves Program Semantics]
\lean{Complexity.LogProg.PS.of_eq}
\label{Complexity.LogProg.PS.of_eq}
\leanok
\uses{Complexity.LogProg.PReg, Complexity.LogProg.PS, Complexity.LogProg.PathLabel}
\difficulty{1}
Let $x$ be a finite Boolean word, $B$ a natural number, $l$ a program label, $\sigma,\sigma'$
register states, and $P$ a proposition. If $\sigma=\sigma'$ and, starting at $l$ with register
state $\sigma'$, every computation terminates with all registers bounded by $B$ at every
preterminal configuration satisfying the program’s safety conditions, and an accepting
computation exists exactly when $P$ holds, then the same properties hold starting at $l$
with register state $\sigma$.
\end{lemma}

\begin{lemma}[Matrix row scan]
\lean{Complexity.LogProg.ps_row}
\label{Complexity.LogProg.ps_row}
\leanok
\uses{Complexity.LogProg.NSem.clr, Complexity.LogProg.NSem.inc, Complexity.LogProg.NSem.jeq_neg, Complexity.LogProg.NSem.jeq_pos, Complexity.LogProg.PReg, Complexity.LogProg.PReg.Le, Complexity.LogProg.PReg.reg, Complexity.LogProg.PS, Complexity.LogProg.PS.of_eq, Complexity.LogProg.nsem_addTwice, Complexity.LogProg.pathA, Complexity.LogProg.pathOr}
\difficulty{6}
Let $x$ be a finite bit string, $B$ a natural number, $P$ a proposition, and
$\sigma=(p,n,a,b,s,t,tmp,st,nx)$ a program register state. Suppose $2|x|+2\le B$, every
register of $\sigma$ is at most $B$, $a\le s<n$, $p=2n+2+2an$, $b\le n$, and $2n+2+2n^2\le
|x|$. Assume that for every $b'\le n$, the program started at the column-check stage
with row register $s$, column register $b'$, and pointer $2n+2+2sn$ terminates with all
registers bounded by $B$ and has an accepting computation exactly when $P$ holds. Then
the program started at the row-scan stage in state $\sigma$ terminates with all registers
bounded by $B$ and has an accepting computation exactly when $P$ holds.
\end{lemma}

\begin{lemma}[Path program phase transfer]
\lean{Complexity.LogProg.ps_index}
\label{Complexity.LogProg.ps_index}
\leanok
\uses{Complexity.LogProg.NSem.clr, Complexity.LogProg.NSem.inc, Complexity.LogProg.PReg, Complexity.LogProg.PReg.Le, Complexity.LogProg.PReg.reg, Complexity.LogProg.PS, Complexity.LogProg.PS.of_eq, Complexity.LogProg.nsem_addTwice, Complexity.LogProg.pathA, Complexity.LogProg.pathOr, Complexity.LogProg.ps_row, Complexity.rowMajor_lt}
\difficulty{5}
Let $x$ be a finite Boolean word, $B$ a natural number, $P$ a proposition, and $\sigma$
a program register state whose nine register values are all at most $B$. Suppose
$2|x|+2\le B$, the source $s$ and guessed next vertex $nx$ satisfy $s<n$ and $nx<n$, and
$2n+2+2n^2\le |x|$, where $n$ is the vertex-count register. If the path program, started
at the label $wrd$ with input pointer $2n+2+2(sn+nx)$, row register $s$, and column
register $nx$, has the termination, register-bound, and acceptance properties expressed
by $PS$ with conclusion $P$, then the same properties hold when it is started at label
$wcp$ with register state $\sigma$ and conclusion $P$.
\end{lemma}

\begin{theorem}[Bounded walk computation]
\lean{Complexity.LogProg.ps_walk}
\label{Complexity.LogProg.ps_walk}
\leanok
\uses{Complexity.LogProg.NSem.call_neg, Complexity.LogProg.NSem.call_pos, Complexity.LogProg.NSem.clr, Complexity.LogProg.NSem.congr, Complexity.LogProg.NSem.inc, Complexity.LogProg.NSem.jeq_neg, Complexity.LogProg.NSem.jeq_pos, Complexity.LogProg.NSem.ret, Complexity.LogProg.PReg, Complexity.LogProg.PReg.Le, Complexity.LogProg.PReg.reg, Complexity.LogProg.PReg.v_0, Complexity.LogProg.PReg.v_1, Complexity.LogProg.PReg.v_8, Complexity.LogProg.PS, Complexity.LogProg.PS.of_eq, Complexity.LogProg.nsem_copy, Complexity.LogProg.nsem_guess, Complexity.LogProg.pathA, Complexity.LogProg.pathOr, Complexity.LogProg.pathOr_bit, Complexity.LogProg.ps_index, Complexity.WalkLe, Complexity.pathEdge, Complexity.rowMajor_lt}
\difficulty{6}
Let $x$ be a finite bit string, let $B$ be a natural number, and let
$\sigma=(p,n,a,b,s,t,tmp,st,nx)$ be a nine-register state. Suppose $2|x|+2\le B$, every
register of $\sigma$ is at most $B$, $s<n$, $t<n$, $st\le n$, and $2n+2+2n^2\le |x|$.
Define a directed edge relation on the natural numbers by declaring $u$ adjacent to $v$
when the bit of $x$ at position $2n+2+2(un+v)$ is true, with out-of-range entries
treated as false. Then the program started at label $w0$ in state $\sigma$ terminates on
every branch with all registers at most $B$, and has an accepting branch exactly when
there is a walk from $s$ to $t$ of length at most $n-st$.
\end{theorem}
